\section{Patterns in the Corpus}

\input{figures/FIG_timeline.tex}
\input{figures/FIG_network.tex}

The tables and figures in this section are deliberately simple. Their purpose is not to turn a small humanities corpus into a false science. It is to make the paper's descriptive claims easier to inspect.

First, the event timeline shows that Frederick-linked records are not evenly distributed. They cluster in two meaningful zones. One is the earlier zone of Maryland footing and reputational management, including the Monocacy context and the 1804 Adams letter. The other is the rehabilitation zone of 1811 and 1812, where Frederick becomes the court-martial site and the place from which Wilkinson appeals directly to presidential fairness. The point is not merely that Frederick appears often. It appears when Wilkinson is most in need of repair.

Second, the weighted connection network now places Frederick Town as the heaviest place node in the corpus and the second-heaviest node overall, behind only Wilkinson himself. That ranking is historically useful. It means Frederick is not just a side label attached to one trial. It connects outward to Adams, Madison, Eustis, Taney, the barracks, and Wilkinson himself. In a small, curated corpus, that kind of centrality is best read as a descriptive pattern rather than a mechanical proof, but it does show that the local place is carrying more than anecdotal weight.

Third, the corpus metrics help separate intimacy from function. Only one record in the full set is primarily a social-access document. By contrast, sixteen turn on use, protection, defense, procedure, testimony, or settlement. This does not prove that Jefferson disliked Wilkinson or that Wilkinson felt no real attachment. It does show that the documentary relationship, as it survives on paper, was mostly organized around use, defense, and political adjudication.

Fourth, the theme counts are modest but suggestive. Wilkinson's coded excerpts contain more defense language than warmth language; Jefferson's contain more utility language than warmth language. The pattern fits the historical reading. Wilkinson wrote upward to save himself. Jefferson wrote downward to steady an officer he still found usable, then later to preserve a memory of service without reopening every scandal.

\begin{table}[t]
\centering
\small
\begin{tabular}{lrrr}
\toprule
Author & Warmth / 1000 tok. & Utility / 1000 tok. & Defense / 1000 tok. \\
\midrule
James Wilkinson & 10.93 & 13.66 & 16.39 \\
Thomas Jefferson & 10.93 & 21.86 & 10.93 \\
\bottomrule
\end{tabular}
\caption{Normalized theme rates per 1000 tokens in the coded excerpts and summaries. The normalized view still leaves Wilkinson's voice more defense-heavy than warmth-heavy, while Jefferson remains most associated with utility language.}
\label{tab:themerates}
\end{table}


The normalized theme-rate check in Table~\ref{tab:themerates} makes the same point more conservatively. Once token volume is taken into account, Wilkinson still runs defense-heavier than warmth, while Jefferson remains most closely associated with utility language. That does not turn a small lexicon into sentiment science. It simply shows that the asymmetry survives a stricter denominator.

Fifth, the evidence-tier counts help show why the local argument is stronger than a ceremonial hometown claim. Nearly three quarters of the Frederick-linked corpus is primary. The Frederick set now includes not just an early reputational-defense letter, a plea for fairness, and an acquittal notice, but also venue planning, jurisdiction management, witness coordination, interrogatory practice, and testimonial return. That balance matters because it means the paper is not deriving its Frederick thesis from later local boosterism alone.

\begin{table}[t]
\centering
\small
\begin{tabular}{lr}
\toprule
Primary-only node & Weight \\
\midrule
James Wilkinson & 32 \\
Thomas Jefferson & 18 \\
Frederick Town & 18 \\
James Madison & 9 \\
Aaron Burr & 7 \\
\bottomrule
\end{tabular}
\caption{Top nodes in the primary-only connection network. Frederick Town remains one of the heaviest nodes even when the later local-context sources are removed.}
\label{tab:primarynetwork}
\end{table}

\begin{table}[t]
\centering
\small
\begin{tabular}{lrr}
\toprule
Place node & Full corpus & Primary-only \\
\midrule
Frederick Town & 10 & 6 \\
Frederick Barracks & 3 & 0 \\
Monocacy & 2 & 0 \\
\bottomrule
\end{tabular}
\caption{Place-node ranking in the connection network. Frederick Town dominates both in the full corpus and in the primary-only subset, reinforcing its role as the paper's central local geography.}
\label{tab:placenodes}
\end{table}


The primary-only network provides a related robustness check. Once the later local-context materials are removed, Frederick Town remains tied with Jefferson as the second-heaviest node in the connection graph, behind only Wilkinson. That is the right result for this paper. It does not mean Frederick outweighs the national actors. It means Frederick remains structurally visible even in the documentary core.

The place-node ranking sharpens the local point further. Frederick Town outranks the Frederick Barracks and Monocacy both in the full corpus and in the primary-only subset. That matters because it shows the argument is not simply about any Maryland geography touching Wilkinson. It is specifically about Frederick as the most active local place in the documentary chain of survival.

The current revision adds one more modest robustness layer. An excerpt-only lexicon pass and a more conservative dictionary preserve the same directional asymmetry between Wilkinson's defensive language and Jefferson's utility language, while leave-one-out perturbations still leave Frederick Town as the leading place node. These checks do not substitute for second-coder validation, but they make it less likely that the main descriptive pattern is a one-pass artifact of this exact phrasing choice.

These results are best understood as corroboration. They do not replace historical judgment, and the corpus is too selective to support extravagant quantitative claims. They do, however, make one thing harder to deny: if a local society wants to know why Frederick matters in the Wilkinson story, the answer is visible not only in prose but in the shape, timing, and evidentiary layering of the surviving record itself.
